\documentclass{beamer}
\usepackage{ctex}
\usetheme{SimplePlus}
\setbeamertemplate{footline}{}

\usepackage{amsmath,amssymb,mathtools}
\usepackage{tikz}
\usetikzlibrary{arrows.meta}
\usepackage{tikz-cd}

\usepackage{unicode-math}
%\setmathfont{STIX Two Math}
%\setmathfont{TeX Gyre Termes Math}
%\setmathfont{Libertinus Math}
%\setmathfont{Fira Math}

\usepackage[
  hybrid,                            % 允许在 Markdown 中混用 LaTeX 命令
  texMathDollars                     % 启用 $...$ 和 $$...$$ 数学公式
]{markdown}

\title{modality 相关直觉}

\date{}

\begin{document}

\begin{frame}
    % Print the title page as the first slide
    \titlepage
\end{frame}

\begin{frame}[fragile]{Adjunction 的基本定义}
\begin{markdown}
设有两个范畴：

$$
\mathcal C,\mathcal D
$$

以及两个函子：

$$
\begin{gathered}
F:\mathcal C\to\mathcal D \\
G:\mathcal D\to\mathcal C
\end{gathered}
$$

如果存在自然同构：

$$
\mathrm{Hom}_{\mathcal D}(F C,D)
\cong
\mathrm{Hom}_{\mathcal C}(C,G D)
$$

则称：

$$
F\dashv G
$$

读作：

> $F$ 是 $G$ 的左伴随，$G$ 是 $F$ 的右伴随。
\end{markdown}
\end{frame}

\begin{frame}[fragile]{为什么需要 unit 和 counit？}
\begin{markdown}
伴随的 Hom-set 同构是一种“外部描述”。

但是范畴论希望把它内部化。

也就是说不用：
$$
\mathrm{Hom}(FC,D)
\cong
\mathrm{Hom}(C,GD)
$$

而直接给出两个自然变换：

$$
\eta
$$

和

$$
\varepsilon
$$

它们叫：

* unit
* counit
\end{markdown}
\end{frame}

\begin{frame}[fragile]{Unit 的定义}
\begin{markdown}
对于：

$$
F\dashv G
$$

unit 是：

$$
\boxed{
\eta:
Id_{\mathcal C}
\Rightarrow
GF
}
$$

也就是一个自然变换：

$$
\eta_C:C\to GF(C)
$$

对于每个对象：

$$
C\in\mathcal C
$$

都有一个态射：

$$
\eta_C:C\rightarrow G(F(C))
$$
\end{markdown}
\end{frame}

\begin{frame}[fragile]{Counit 的定义}
\begin{markdown}
counit 是：

$$
\boxed{
\varepsilon:
FG
\Rightarrow
Id_{\mathcal D}
}
$$

也就是：

对于：

$$
D\in\mathcal D
$$

有：

$$
\varepsilon_D:
F(GD)\rightarrow D
$$
\end{markdown}
\end{frame}

\begin{frame}[fragile,shrink]{Unit 和 counit 的关系\cite{ncatlab_triangle_identities}}
\begin{markdown}
它们不是任意自然变换。必须满足两个三角恒等式：

---

第一个 triangle identity
$$
F
\xrightarrow{F\eta}
FGF
\xrightarrow{\varepsilon F}
F
$$

必须等于：

$$
id_F
$$

写成：

$$
\boxed{
\varepsilon_F\circ F\eta=id_F
}
$$

展开：

$$
F C
\rightarrow
FGFC
\rightarrow
FC
$$

结果：

$$
FC
$$

不变。
\end{markdown}
\end{frame}

\begin{frame}[fragile]
\begin{markdown}
第二个 triangle identity
$$
G
\xrightarrow{\eta G}
GFG
\xrightarrow{G\varepsilon}
G
$$

必须：

$$
\boxed{
G\varepsilon\circ\eta G=id_G
}
$$
\end{markdown}
\end{frame}

\begin{frame}[fragile,shrink]{和 Monad / Comonad 的关系}
\begin{markdown}
如果：

$$
F\dashv G
$$

那么：

$$
GF
$$

自动形成一个 Monad：

$$
T=GF
$$

其中：

unit：

$$
\eta:
Id\Rightarrow GF
$$

就是 Monad 的 unit。

multiplication：

$$
\mu:
GFGF
\rightarrow
GF
$$

定义：

$$
G\varepsilon F
$$

即：

$$
GFGF
\xrightarrow{G\varepsilon F}
GF
$$
\end{markdown}
\end{frame}

\begin{frame}[fragile]
同样：

$$
FG
$$

形成 Comonad：

$$
\varepsilon:
FG\Rightarrow Id
$$

就是 counit。
\end{frame}

\begin{frame}[fragile]{reflective subcategory}
\begin{markdown}
设
$$
\mathcal D \xhookrightarrow{i} \mathcal C
$$

是一个全忠实子范畴。如果包含函子 $i$ 有左伴随

$$
L:\mathcal C\rightleftarrows\mathcal D:i,
\qquad L\dashv i,
$$

那么 $\mathcal D$ 称为 $\mathcal C$ 的 **reflective subcategory**，而

$$
L(X)
$$

称为 $X$ 在 $\mathcal D$ 中的 **reflection**。

伴随的 unit 给出一个自然映射：

$$
\eta_X:X\to iL(X).
$$

这就是把 $X$ “反射”到子范畴 $\mathcal D$ 中。
\end{markdown}
\end{frame}

\begin{frame}[fragile]

它满足普适性质：对于任意 $D\in\mathcal D$，所有从 $X$ 到 $iD$ 的映射都唯一地经过 $\eta_X$ 分解：

\[
\begin{tikzcd}
X \arrow[r,"\eta_X"] \arrow[dr,"f"'] & iL(X) \arrow[d,"i(g)"] \\
& iD
\end{tikzcd}
\]

也就是说，对任意
$$
f:X\to iD
$$

存在唯一的
$$
g:L(X)\to D
$$
使得
$$
f=i(g)\circ\eta_X.
$$

所以 $L(X)$ 是 $X$ 在 $\mathcal D$ 里的“最通用近似”。
\end{frame}

\begin{frame}[fragile,shrink]{一个典型例子：Ab 是 Grp 的 reflective subcategory}
\begin{markdown}

令

$$
\mathbf{Ab}\hookrightarrow\mathbf{Grp}
$$

是交换群包含到群中的子范畴。

每个群 $G$ 都有它的交换化：

$$
G_{\mathrm{ab}}=G/[G,G].
$$

有自然商映射：

$$
\eta_G:G\to G_{\mathrm{ab}}.
$$

对任何交换群 $A$，任意群同态

$$
f:G\to A
$$

都会唯一地经过 $G_{\mathrm{ab}}$：

$$
G\to G_{\mathrm{ab}}\to A.
$$

因此：

$$
(-)_{\mathrm{ab}}:\mathbf{Grp}\to\mathbf{Ab}
$$

是包含函子的左伴随，所以 $\mathbf{Ab}$ 是 $\mathbf{Grp}$ 的 reflective subcategory。

这里“反射”的直观意思就是：

> 把任意群 $G$ 映射到一个最通用的交换群 $G_{\mathrm{ab}}$。

而且如果 $G$ 本来已经是交换群，那么：

$$
G_{\mathrm{ab}}\cong G.
$$

这正对应“反射两次没有新变化”，即幂等 monad：

$$
L i\cong \mathrm{id}_{\mathbf{Ab}}.
$$

对偶地，如果包含函子有右伴随，就叫 **coreflective subcategory**，因为那对应的是从一般对象中“取出”子范畴中的最大或最通用部分，而不是把它映射进去。
\end{markdown}
\end{frame}

\begin{frame}[fragile,shrink]{coreflective subcategory 的典型例子}
\begin{markdown}
一个非常典型的例子是：

$$
\mathbf{TorAb}\hookrightarrow \mathbf{Ab},
$$

即**扭群（torsion abelian groups）**包含到所有阿贝尔群中。

对任意阿贝尔群 $A$，定义其扭部分：

$$
t(A)=\{a\in A\mid \exists n>0,\;na=0\}.
$$

显然 $t(A)$ 是一个扭群，并且有自然包含映射：

$$
\varepsilon_A:t(A)\hookrightarrow A.
$$

这个映射具有如下普适性质：对于任意扭群 $T$，每个同态

$$
f:T\to A
$$

都会唯一地经过 $t(A)$ 分解：
\end{markdown}
\[
\begin{tikzcd}
T \arrow[r,"f"] \arrow[dr,"\bar{f}"'] & A \\
& t(A) \arrow[u,"\epsilon_A"']
\end{tikzcd}
\]
\begin{markdown}
也就是存在唯一的同态

$$
\bar f:T\to t(A)
$$

使得

$$
f=\varepsilon_A\circ\bar f.
$$
\end{markdown}
\end{frame}

\begin{frame}[fragile,shrink]
\begin{markdown}
原因是：如果 $x\in T$ 是扭元素，那么 $f(x)$ 也是扭元素。因此 $f(T)\subseteq t(A)$。

所以包含函子

$$
i:\mathbf{TorAb}\hookrightarrow\mathbf{Ab}
$$

有右伴随

$$
t:\mathbf{Ab}\to\mathbf{TorAb},
$$

即

$$
i\dashv t.
$$

因此 $\mathbf{TorAb}$ 是 $\mathbf{Ab}$ 的一个 **coreflective subcategory**。

这里的 coreflection 是：

$$
t(A)\hookrightarrow A.
$$

它表示：

> 从任意阿贝尔群 $A$ 中提取出其中最大的扭部分 $t(A)$。

并且如果 $A$ 本身已经是扭群，那么

$$
t(A)\cong A.
$$

因此这个 coreflection 是幂等的：

$$
t(t(A))\cong t(A).
$$

相应的 endofunctor

$$
G=i\circ t:\mathbf{Ab}\to\mathbf{Ab}
$$

是一个 **idempotent comonad**。它的：

- counit 是包含映射

  $$
  \varepsilon_A:t(A)\hookrightarrow A;
  $$

- comultiplication 是

  $$
  \delta_A:t(A)\to t(t(A)),
  $$

  而这个映射实际上是同构，因为

  $$
  t(t(A))=t(A).
  $$

和 reflective subcategory 对比：

$$
\begin{array}{ccl}
\text{reflective} &:& X\to L(X)\\
\text{coreflective} &:& R(X)\to X
\end{array}
$$

在这个例子中：

$$
\text{反射：}\quad A\to A_{\mathrm{ab}}
$$

可以把群变成最通用的交换群；而

$$
\text{余反射：}\quad t(A)\hookrightarrow A
$$

则是从阿贝尔群中提取最大的扭子群。
\end{markdown}
\end{frame}

\begin{frame}[fragile,shrink]{Monad 与 idempotent monad}
\begin{markdown}
一个 monad 是三元组

$$
(T,\eta,\mu)
$$

其中：

$$
T:\mathcal C\to\mathcal C
$$

$$
\eta_X:X\to TX
$$

是 unit，

$$
\mu_X:T^2X\to TX
$$

是 multiplication，并满足结合律和单位律。

如果 $\mu_X$ 对每个 $X$ 都是同构：

$$
T^2X \simeq TX,
$$

那么这个 monad 称为 **idempotent monad**，也叫 **Kock–Zöberlein monad** 的某些特殊情形，最常见的范畴论表述是：

> 幂等 monad 对应于一个反射子范畴（reflective subcategory）。

设

$$
L:\mathcal C\rightleftarrows \mathcal D:i
$$

是一个伴随，其中 $i$ 是全忠实函子：

$$
L\dashv i.
$$

那么复合

$$
T=iL
$$

就是一个幂等 monad。这里的 $TX$ 可以理解为把 $X$ “反射”到某个特殊子范畴中的对象。

例如：

- 集合 $\to$ 群的自由群 monad 不是幂等的，因为自由群再做一次并不会什么都不改变；
- 集合 $\to$ 命题截断 $\|X\|$ 在 HoTT 中是幂等的，因为

  $$
  \|\|X\|\|\simeq \|X\|.
  $$
\end{markdown}
\end{frame}

\begin{frame}[fragile,shrink]{Comonad 与 idempotent comonad}
\begin{markdown}
comonad 是 monad 的对偶。它由

$$
(G,\varepsilon,\delta)
$$

组成，其中：

$$
G:\mathcal C\to\mathcal C
$$

$$
\varepsilon_X:GX\to X
$$

是 counit，

$$
\delta_X:GX\to G^2X
$$

是 comultiplication。

如果

$$
\delta_X:GX\to G^2X
$$

是同构，即

$$
GX\simeq G^2X,
$$

那么 $G$ 是 **idempotent comonad**。

它对应于一个 **余反射子范畴（coreflective subcategory）**。若

$$
i:\mathcal D\rightleftarrows\mathcal C:R
$$

满足

$$
i\dashv R
$$

且 $i$ 全忠实，则

$$
G=iR
$$

是一个幂等 comonad。

直观地说：

- idempotent monad：把对象“自由地变成某种好对象”；
- idempotent comonad：从对象中“提取出某种满足条件的部分”。
\end{markdown}
\end{frame}

\begin{frame}[fragile,shrink]{Modality 是什么？}
\begin{markdown}
在现代类型论，尤其是 HoTT 中，**modality** 通常指一个特殊的 endofunctor

$$
\bigcirc:\mathcal U\to\mathcal U
$$

并带有自然变换

$$
\eta_X:X\to \bigcirc X.
$$

它满足

$$
\bigcirc\bigcirc X\simeq \bigcirc X.
$$

所以从范畴论角度看，一个 modality 通常是：

$$
\boxed{\text{modality}=\text{idempotent monad}}
$$

也就是说，modality 的结构通常包括：

$$
\eta:X\to \bigcirc X
$$

以及

$$
\mu:\bigcirc\bigcirc X\to \bigcirc X,
$$

并且 $\mu$ 是等价。

因此有：

$$
\bigcirc^2 X\simeq \bigcirc X.
$$

其固定对象满足：

$$
X\simeq \bigcirc X.
$$

这些固定对象形成一个反射子范畴，通常称为 **modal types**。
\end{markdown}
\end{frame}

\begin{frame}[fragile,shrink]{为什么 modality 经常还要求 left exact？}
\begin{markdown}
在 HoTT 中，常见的 modality 不仅要求幂等 monad，还要求它保持有限极限，特别是保持 dependent sums 或 pullback。这样的 modality 称为 **left exact modality**。

粗略地说：

$$
\text{left exact modality} = \text{idempotent monad} + \text{preservation of finite limits}.
$$

left exact 条件使 modality 能够良好地作用于：

- dependent types；
- identity types；
- pullbacks；
- proposition；
- sheaf 和 stack 等结构。

例如：

- propositional truncation 是一种 modality；
- $n$-truncation 是一种 modality；
- sheafification 通常给出一种 left exact modality；
- localization 也常由 modality 描述。

所以要注意：

> 不是每一个幂等 monad 都自动满足 HoTT 中通常要求的全部 modality 条件；是否称为 modality，取决于所采用的定义。
\end{markdown}
\end{frame}

\begin{frame}[fragile,shrink]{ Modality 和 comonad 的关系}
\begin{markdown}
如果一个 modality 是 monad，那么它有一个单位映射：

$$
\eta_X:X\to \bigcirc X.
$$

这表示把 $X$ 送到它的“模态化版本”。

例如命题截断：

$$
X\to \|X\|.
$$

这是从一般对象到更“粗糙”或更“局部化”的对象。

其对偶操作则是 comonad：

$$
\Box X\to X.
$$

这通常表示从 $X$ 中提取出“满足某种模态条件的部分”。因此：

$$
\boxed{
\text{monadic modality：} X\to \bigcirc X
}
$$

$$
\boxed{
\text{comonadic modality：} \Box X\to X
}
$$

如果 $\bigcirc$ 是幂等 monad，那么它的固定对象是 reflective subcategory；如果 $\Box$ 是幂等 comonad，那么它的固定对象是 coreflective subcategory。
\end{markdown}
\end{frame}

\begin{frame}[fragile,shrink]{逻辑中的直观区别}
\begin{markdown}
可以把两者类比成逻辑中的两个方向。

### Monad modality

$$
X\to \bigcirc X
$$

表示：

> 把 $X$ 转换成满足某种性质的版本。

例如：

$$
X\to \|X\|
$$

把一个类型变成命题。

或者：

$$
X\to \|X\|_n
$$

把它变成 $n$-截断类型。

### Comonad modality

$$
\Box X\to X
$$

表示：

> 从 $X$ 中取出一个更受限制、更结构化的部分。

在模态逻辑中，常见的是：

$$
\Box A\to A
$$

但要注意，普通模态逻辑中的 $\Box$ 不一定是幂等的。只有当

$$
\Box\Box A\leftrightarrow \Box A
$$

成立时，才对应幂等 comonad。

例如：

- S4 的 $\Box$ 满足通常的幂等性条件；
- K、T 等一般不一定满足幂等性。
\end{markdown}
\end{frame}

\begin{frame}[fragile,shrink]{Monad 和 comonad 是否互相转换？}
\begin{markdown}
单独给一个 monad，一般不能自动得到一个 comonad。

但是，如果有伴随

$$
L\dashv R,
$$

那么会同时产生：

$$
RL
$$

上的 monad，以及

$$
LR
$$

上的 comonad：

$$
\mathcal C
\overset{L}{\longrightarrow}
\mathcal D
\overset{R}{\longrightarrow}
\mathcal C.
$$

具体地：

- $RL$ 是 $\mathcal C$ 上的 monad；
- $LR$ 是 $\mathcal D$ 上的 comonad。

如果其中的某个函子全忠实，就得到幂等性：

- $R$ 全忠实 $\Rightarrow RL$ 是幂等 monad；
- $L$ 全忠实 $\Rightarrow LR$ 是幂等 comonad。
\end{markdown}
\end{frame}

\begin{frame}[fragile,shrink]{最简关系图}
\begin{markdown}
$$
\begin{array}{ccc}
\text{reflective subcategory}
&\longleftrightarrow&
\text{idempotent monad}
\\
\updownarrow && \updownarrow
\\
\text{modal types}
&\longleftrightarrow&
\text{modality}
\end{array}
$$

对偶地：

$$
\begin{array}{ccc}
\text{coreflective subcategory}
&\longleftrightarrow&
\text{idempotent comonad}
\\
\updownarrow && \updownarrow
\\
\text{comodal/coreflective types}
&\longleftrightarrow&
\text{comonadic modality}
\end{array}
$$

一句话总结：

$$
\boxed{
\text{HoTT 中通常的 modality 是 left exact 的 idempotent monad；}
}
$$

$$
\boxed{
\text{其对偶概念是 idempotent comonad，对应 coreflective subcategory。}
}
$$

但“modality”这个词有时专指 monadic 方向；comonad 方向通常会明确称作 **comonadic modality**、**comodality**，或者直接说 **idempotent comonad**。
\end{markdown}
\end{frame}

% 参考文献帧，长列表自动分页
\begin{frame}[allowframebreaks]{参考文献}
\begin{thebibliography}{99}

\bibitem{ncatlab_triangle_identities}
\url{https://ncatlab.org/nlab/show/triangle+identities}

\end{thebibliography}
\end{frame}
\end{document}